\section{Related work and research gap}
\label{sec:related}

\subsection{Industrial anomaly detection and process graphs}

Industrial-control anomaly detection is usually evaluated on multivariate telemetry whose variables are coupled by control logic and physical process relationships. The HAI testbed combines physical control systems with hardware-in-the-loop simulation and provides normal and attack traces suitable for chronological evaluation \citep{shin2020hai,shin2021two}. Evaluation is itself nontrivial: point-wise accuracy can reward detectors that find only a small portion of a long episode, motivating time-series-aware measures such as eTaPR \citep{hwang2022etapr} and broader benchmark discipline \citep{schmidl2022evaluation}.

Graph and attention models provide a natural representation of interacting sensors and actuators. Graph attention aggregates information over related variables \citep{velickovic2018gat}, while Transformer attention captures long-range temporal dependence \citep{vaswani2017attention}. Graph-Transformer anomaly detectors have been studied for multivariate IoT series \citep{chen2022gta} and controller-area networks \citep{zhou2025cgts}; a recent survey organizes the wider graph-anomaly landscape \citep{qiao2025deep}. Our detector follows this engineering pattern but keeps relation provenance separate: official process relations, device-channel links, operational context, normal-training dependencies, and constrained hardware-in-the-loop couplings are not presented as equally authoritative physical edges.

The detector is a means rather than the paper's sole contribution. We report window, episode, localization, false-alarm, and calibration outcomes because explanation quality cannot be separated from the evidence model that produced the explanation. In particular, a perfectly enforced explanation contract cannot repair a missed attack or a poorly calibrated alarm.

\subsection{Explainable and human-centered cybersecurity}

Surveys of explainable cybersecurity emphasize transparency, intelligence automation, and operator trust \citep{zhang2022xaicyber,sarker2024xai,sharma2025review}. Human-centered systems increasingly combine detection with understandable narratives, and recent work proposes multi-agent CPS security workflows \citep{ferraro2026multiagent} and human-centric explanation frameworks \citep{barik2026deepguard}. LLM-based studies generate cyberattack explanations for automatic-generation-control systems \citep{sharshar2025llm} or use prompt-configured models as complementary Modbus triage layers \citep{kong2026llm}. LLMs have also been used offensively to generate industrial attack patterns \citep{ahmed2025attackllm}. Together, these studies show that generative models can operate on both sides of the security boundary.

Human factors prevent a simple equation of ``more explanation'' with ``better decision.'' Explanations can raise reliance even when the recommendation is wrong \citep{bansal2021whole}; cognitive forcing can reduce overreliance \citep{bucinca2021trust}; and security analysts face domain-specific automation and confirmation biases \citep{hagen2025human}. These findings motivate an operator-facing contract that controls what the model is allowed to say, not only whether its words sound plausible.

\subsection{Adversarial explanations}

Post-hoc explanation methods can be manipulated independently of predictive behavior. \citet{slack2020fooling} demonstrated that LIME and SHAP can be made to conceal undesirable classifier behavior. Cybersecurity studies subsequently analyzed black-box and adversarial XAI attacks \citep{kuppa2020blackbox,kuppa2021adversarial}. Recent surveys broaden the threat taxonomy to model, data, backdoor, and human-facing attacks \citep{baniecki2024survey,vadillo2025survey}; 2025 experiments continue to find vulnerabilities in cybersecurity explanation pipelines \citep{mia2025vulnerable}. These risks align with the broader NIST adversarial-ML taxonomy \citep{vassilev2025aml} and generative-AI risk-management guidance \citep{nist2024genai}.

Our threat differs in two ways. First, the attacker does not need to alter the detector input, prediction, or committed attribution set. Second, success is defined at the deployment gate: a materially changed operator surface passes the verifier. This makes the study narrower than human deception research, because we do not infer actual harmful influence from text alone, but more operational than attribution-only robustness because the attacked object is the message shown to an industrial operator.

\subsection{Gap and design response}

The literature establishes that explanations can be useful, over-trusted, and attacked. It does not yet provide a reproducible answer to the following engineering gap: how can an industrial AI system allow generative assistance while making a closed set of decision-critical statements non-negotiable? \system{} answers by moving those statements out of the generative trust boundary. Rather than asking a semantic classifier to recognize every possible manipulative paraphrase, the system renders integrity-critical fields canonically and verifies equality. A finite registered extension permits only prereviewed benign renderings. The benefit is a simple, auditable guarantee within the contract. The cost is reduced linguistic flexibility, which we measure and discuss rather than hiding.
